\subsection{三捩れtranslationの量子Weyl chartと全中心}
\label{subsec:quantum3-center}
\begin{proposition}[三捩れchartの明示的な全中心]
\label{prop:quantum3-full-center}
$q^2+q+1=0$ とし、
\[
 V= k\langle x,y,z\rangle/(yz-qzy-b_1,\ zx-qxz-b_2,\ xy-qyx-b_3)
\]
と置く。$d=(q-1)^{-1}=-(q+2)/3$、$e=-(2q+1)/9$ とすると、
\[
 X=x^3,\quad Y=y^3,\quad Z=z^3,\quad
 W=xyz+d(b_1x+qb_2y+b_3z)
\]
は中心元であり、全中心は
\[
 Z(V)=k[X,Y,Z,W]/
 \bigl(W^3-db_1b_2b_3W-XYZ-e(b_1^3X+b_2^3Y+b_3^3Z)\bigr).
\]
$\lambda^3\ne1$ のとき、中心切断
\[
 U_D=V/(X+Y+Z-3\lambda W+D)
\]
のassociated gradedは滑らかな楕円曲線の次数3 THCRで、translationは正確に位数3を持つ。
さらに $Z(U_D)=Z(V)/(X+Y+Z-3\lambda W+D)$ である。
従って全中心は表示済みの二方程式で与えられ、Rees環のH3性は中心McKay判定と同値である。
\end{proposition}
\begin{proof}
順序 $x<y<z$ に対する書換えは
$yx=q^2xy-q^2b_3$、$zx=qxz+b_2$、$zy=q^2yz-q^2b_1$ である。
唯一の重なり $zyx$ の二還元は一致し、PBW基底は $x^iy^jz^k$ となる。
上記四元の全生成元との交換子と表示された次数9恒等式は、
$q^2+q+1$ で係数を還元して直接検証できる。

最高次数の量子多項式環で、単項式 $x^iy^jz^k$ が中心である条件は
$i\equiv j\equiv k\pmod3$ である。
従ってその中心は $x^3,y^3,z^3,xyz$ で生成される。
これら全てが $V$ の中心へ持ち上がっているから、
最高項を逐次引くことにより全中心の生成を得る。
関係を $W$ に関するmonic三次で割った剰余の最高項は、
$X^aY^bZ^cW^r$、$r=0,1,2$ の相異なるPBW単項式になる。
従って他の関係は存在しない。

$S=\operatorname{gr}V$ の中心三次
$g=x^3+y^3+z^3-3\lambda xyz$ は滑らかなHesse曲線を定める。
三二次関係のgraphは $(x:y:z)\mapsto(x:qy:q^2z)$ であり、
この曲線上では非零三捩れ点によるtranslationとなる。
section環への写像は三二次関係と $g$ を殺し、
Hilbert級数の一致により $S/(g)$ とTHCRは一致する。
$g$ は整域 $S$ で正則なので、中央filtered元による商の最高記号は
正確に $(g)$ となり、$\operatorname{gr}U_D=S/(g)$ を得る。
中心の商との比較も証明する。$f=X+Y+Z-3\lambda W+D$ と置く。
$fv$ が中心なら $f[v,a]=0$ と整域性から $v$ も中心だから、
$Z(V)/(f)\hookrightarrow Z(U_D)$ である。
$Z(S)/(g)\hookrightarrow Z(S/(g))$ も同じ議論で単射となる。
前者のHilbert級数は $(1-s^3)/(1-s)^3$、ここで $s$ は次数3を表す。
後者は $E/\langle\rho\rangle$ 上の次数3 section環なので同じ級数を持ち、両者は一致する。
Rainsのfiltered中心定理により $\operatorname{gr}Z(U_D)=Z(S/(g))$ だから、
最高項の帰納法で $Z(V)/(f)\to Z(U_D)$ は全射でもある。
定理\ref{thm:all-filtered-elliptic-strong-envelope}と
有限位数の認識・中心McKay判定を適用する。
\end{proof}

このchartでH3となるための局所条件は具体的には次の通りである。
滑らかな中心点では自動的にAzumayaである。
$A_1$ 点では二つの単純加群の次元が $(1,2)$、
$A_2$ 点では三つの相異なる一次元単純加群が必要十分である。
他の特異型はPI次数3の正則McKay表現を収容できず除外される。
これは定理\ref{thm:finite-order-elliptic-admissibility}の
$\sum r_\chi\dim\chi=3$ から従う。
全 $b_i,\lambda,D$ の明示的な許容条件の消去と、
各中心点の単純加群の分配の計算は残る。

\begin{theorem}[固定三捩れ包絡の全切断判定]
\label{thm:quantum3-fixed-all-cuts}
前命題で $b_1=b_2=b_3=1,\lambda=0$ とする。
$e=-(2q+1)/9$ と置くと、$G_D=x^3+y^3+z^3+Dt^3$ による
全切断のうち強いH3となる必要十分条件は
\[
 (D^2+4e)(D^2+36e)\ne0
\]
である。$D^2+9e=0$ の二値では中心が $A_1$ 特異点を持つが、切断はH3である。
従って非H3値は正確に四つである。
\end{theorem}
\begin{proof}
中心式から $Z=-D-X-Y$ を消去すると
\[
 F=W^3-dW+XY(X+Y+D)+eD=0,\qquad d^3=e,\quad e^2=-1/27
\]
を得る。臨界条件は
$3W^2=d$、$Y(2X+Y+D)=X(X+2Y+D)=0$ である。
$D\ne0$ では $(X,Y)$ は
$(0,0),(0,-D),(-D,0),(-D/3,-D/3)$ の四点に限る。
初めの三点では $F=0$ の条件を消去すると $D^2+4e=0$、
最後の一点では
\[
 (D^3+27eD)^2-108e=(D^2+9e)^2(D^2+36e)=0
\]
となる。$D=0$ は $W^2=d/3$ と $F=0$ を両立できず特異値でない。
六つの値は相異なり非零である。
各点のHessianは非退化である。実際、$W$ 方向は $6W\ne0$、
$(X,Y)$ ブロックの行列式は初めの三点で $-D^2$、最後で $D^2/3$ である。
従って初めの二値は各三つの $A_1$ 点を、残り四値は各一つの $A_1$ 点を持つ。

一次元表現の方程式は $yz=zx=xy=-d$ なので、
全て $x=y=z=s$、$s^2=-d$ に限る。
中央切断は $D=-3s^3$ であり、対応する二つの値は $D^2+9e=0$ である。
各値で一次元表現は正確に一つ、その中心点は
$X=Y=Z=-D/3$ である。
この $A_1$ 点ではPI次数3と正の多重度により
McKay次元は $(1,2)$ となり、大域次元2を得る。
他の四特異値では一次元表現が存在しない。
$A_1$ 上の有限大域次元PI次数3対称orderには $(1,2)$ が必要だから、これらは非H3である。
非特異中心の値はAzumayaなのでH3であり、全値の判定を尽くす。
\end{proof}

\begin{theorem}[量子Weyl chart全係数の有限admissibility判定]
\label{thm:quantum3-all-parameters}
命題\ref{prop:quantum3-full-center}の任意の $b_1,b_2,b_3,D$ と
$\lambda^3\ne1$ に対し、$Z=-X-Y+3\lambda W-D$ を代入した中心三次を $F(X,Y,W)$ とする。
切断がH3となる必要十分条件は、$F$ の全特異点が次のいずれかを満たすことである。
\begin{enumerate}
\item Hessianの階数が3で、その中心点上に一次元表現が存在する。
\item Hessianの階数が2で、核方向の三次項が非零であり、
その中心点上に相異なる一次元表現が三つ存在する。
\end{enumerate}
一次元表現とその中心点は、次の有限方程式で全て計算できる。
\[
 yz=-db_1,\quad zx=-db_2,\quad xy=-db_3,\qquad
 x^3+y^3+z^3-3\lambda w+D=0,
\]
\[
 (X,Y,W)=(x^3,y^3,w),\qquad
 w=xyz+d(b_1x+qb_2y+b_3z).
\]
存在と個数は幾何学的点で数え、非被約長さを使わない。
このH3部分は係数空間 $\{\lambda^3\ne1\}\times\mathbb A^4_{b_1,b_2,b_3,D}$ の開集合であり、
既出の相対対角分解と任意底変更・fpqc降下に適合する。
\end{theorem}
\begin{proof}
中心は滑らかな楕円最高項を持つ三次曲面なので、特異点集合は有限である。
Hessian階数3は $A_1$、階数2かつ核方向の三次非零は $A_2$ の正確な判定である。
PI次数3で有限大域次元を許すDu Val型はこれらだけである。
$A_1$ の多重度は正の整数の和3なので $(1,2)$、
$A_2$ では三つの正の整数の和3なので $(1,1,1)$ に限る。
他の $A_n$ は少なくとも四指標を持ち、$D,E$ 型の全指標の次元和も3を超える。
$A_1$ で一次元表現があれば、残りのMcKay多重度は2となり全指標が現れる。
$A_2$ では相異なる三つの一次元表現の存在が全指標の出現と同値である。
従って有限位数の全McKay判定を、表示した有限方程式に置き換えられる。

係数空間上のPBW基底は自由であり、最高項の $g$ が正則だから中央切断も次数ごとに平坦である。
境界は全て同じ正確な位数3のtranslationを持つ。
係数空間は被約有限型なので、
定理\ref{thm:finite-order-relative-H3-open}を適用して最後の主張を得る。
\end{proof}

一次元表現の列挙は特に小さい。
三つの $b_i$ が非零なら、切断前に正確に二つの候補を持つ。
正確に二つが非零なら候補は存在しない。
正確に一つが非零なら、例えば $b_1\ne0$ の場合、
$x=0$、$yz=-db_1$、$y^6+Dy^3+(-db_1)^3=0$ に還元される。
全て零なら候補は座標軸上の $x^3+y^3+z^3+D=0$ の点である。
従ってこの定理は未知の高次元単純表現の探索を必要としない。
これは指定した線型項なしの量子Weyl包絡部分の完全な許容判定であり、
三捩れtranslationを持つ全filtered deformationの被覆や全H3分類とは区別する。
